\documentclass[11pt,a4paper]{article}
\usepackage[ngerman]{babel}
\usepackage[T1]{fontenc}
\usepackage[utf8]{inputenc}
\usepackage{amsmath,amssymb}
\usepackage[most]{tcolorbox}
\usepackage{geometry}
\usepackage{enumitem}
\geometry{margin=2.4cm}
\setlength{\parskip}{0.45em}
\setlength{\parindent}{0pt}

\title{Moral, Naturalismus und die Existenz Gottes}
\author{}
\date{}

\begin{document}
	\maketitle
	
	\begin{abstract}
		Dieser Text untersucht zwei miteinander verbundene philosophische Argumente: das moralische Argument gegen den evolutionären Naturalismus und das moralische Argument für die Existenz Gottes. Ausgangspunkt ist die Frage, ob objektive moralische Werte existieren und wie sie sich erklären lassen. Während evolutionäre Prozesse die Entstehung menschlicher Moralvorstellungen erklären können, bleibt umstritten, ob sie auch deren objektive Gültigkeit begründen. Der Theismus bietet hierfür eine mögliche Erklärung. Die Überzeugungskraft der Argumente hängt jedoch von mehreren umstrittenen Voraussetzungen ab, insbesondere vom moralischen Realismus und von der Annahme, dass objektive Moral eine göttliche Grundlage benötigt.
	\end{abstract}
	
	\section{Einleitung}
	
	Moralische Urteile spielen eine zentrale Rolle im menschlichen Leben. Wir beurteilen Handlungen als richtig oder falsch und gehen häufig davon aus, dass bestimmte Handlungen unabhängig von persönlichen Vorlieben moralisch falsch sind. So scheint beispielsweise die ungerechtfertigte Schädigung unschuldiger Menschen nicht bloss eine Frage des individuellen Geschmacks zu sein.
	
	Daraus ergibt sich eine grundlegende philosophische Frage: Existieren moralische Tatsachen objektiv, oder beruhen moralische Urteile letztlich auf menschlichen Einstellungen und gesellschaftlichen Konventionen?
	
	Das moralische Argument für den Theismus setzt bei dieser Frage an. Es behauptet, dass objektive moralische Werte mit einer naturalistischen Weltanschauung schwer zu vereinbaren seien, während ihre Existenz durch die Annahme Gottes besser erklärt werden könne. Dabei sind zwei Argumentationsschritte zu unterscheiden: Zunächst wird die Vereinbarkeit objektiver Moral mit dem Naturalismus infrage gestellt. Anschliessend wird argumentiert, dass die Existenz objektiver Moral für den Theismus spricht.
	
	\section{Zentrale Begriffe}
	
	\textbf{Naturalismus} bezeichnet hier eine Weltanschauung, nach der die Wirklichkeit grundsätzlich durch natürliche Prozesse erklärt werden kann. Der \textbf{evolutionäre Naturalismus} verbindet diese Sichtweise mit der Annahme, dass sich Lebewesen und ihre Eigenschaften durch Evolution entwickelt haben.
	
	\textbf{Moralischer Realismus} ist die Auffassung, dass moralische Aussagen objektiv wahr oder falsch sein können. Moralische Tatsachen wären demnach nicht allein davon abhängig, was einzelne Menschen oder Gesellschaften für richtig halten.
	
	Demgegenüber bestreiten \textbf{moralischer Antirealismus} und verwandte Positionen, dass moralische Urteile auf diese Weise objektive Tatsachen beschreiben. Moralische Aussagen könnten beispielsweise Ausdruck von Gefühlen, Einstellungen oder gesellschaftlichen Regeln sein.
	
	Der \textbf{Theismus} bezeichnet die Auffassung, dass Gott existiert. Im hier untersuchten Argument wird Gott als mögliche Grundlage objektiver moralischer Werte verstanden.
	
	\section{Das moralische Argument gegen den Naturalismus}
	
	Das erste Argument richtet sich gegen die Annahme, dass sich menschliche Moral vollständig durch evolutionäre Prozesse erklären lässt.
	
	Aus evolutionärer Perspektive können moralische Verhaltensweisen einen Überlebensvorteil bieten. Kooperation, gegenseitige Unterstützung und die Einhaltung sozialer Regeln können beispielsweise dazu beitragen, dass Gruppen stabil bleiben und ihre Mitglieder bessere Überlebenschancen haben. Natürliche Selektion kann deshalb erklären, weshalb Menschen bestimmte moralische Intuitionen und Verhaltensweisen entwickelt haben.
	
	Allerdings muss zwischen der Entstehung einer Überzeugung und ihrer Wahrheit unterschieden werden. Eine evolutionäre Erklärung zeigt, weshalb Menschen eine bestimmte moralische Überzeugung besitzen könnten. Sie zeigt für sich genommen nicht, ob diese Überzeugung objektiv richtig ist.
	
	Wenn unsere moralischen Überzeugungen hauptsächlich deshalb entstanden sind, weil sie das Überleben und die Fortpflanzung begünstigten, stellt sich die Frage, weshalb wir davon ausgehen sollten, dass sie objektive moralische Tatsachen zuverlässig erfassen. Der evolutionäre Ursprung einer Überzeugung garantiert ihre Wahrheit nicht.
	
	\begin{tcolorbox}[title=Argument gegen den evolutionären Naturalismus]
		\begin{enumerate}[leftmargin=*]
			\item Wenn der evolutionäre Naturalismus wahr ist, lassen sich unsere moralischen Überzeugungen durch natürliche und evolutionäre Prozesse erklären.
			\item Diese Prozesse können erklären, wie solche Überzeugungen entstanden sind, ohne damit bereits ihre objektive Wahrheit zu begründen.
			\item Wenn objektive moralische Tatsachen existieren, stellt sich deshalb die Frage, wie sie innerhalb einer rein naturalistischen Weltanschauung erklärt werden können.
		\end{enumerate}
	\end{tcolorbox}
	
	Dieses Argument zeigt zunächst eine mögliche Erklärungslücke auf. Es beweist nicht unmittelbar, dass der Naturalismus falsch ist. Insbesondere folgt aus der evolutionären Entstehung moralischer Überzeugungen nicht automatisch, dass diese falsch oder bloss subjektiv sind. Dafür wäre ein zusätzlicher Schritt in der Argumentation erforderlich.
	
	\section{Das moralische Argument für den Theismus}
	
	Das zweite Argument geht von der Möglichkeit aus, dass objektive moralische Werte tatsächlich existieren. Wenn moralische Tatsachen unabhängig von individuellen Meinungen gelten, stellt sich die Frage nach ihrer Grundlage.
	
	Der Theismus bietet eine mögliche Antwort: Gott könnte die Grundlage objektiver moralischer Werte bilden. Wenn es objektive moralische Verpflichtungen gibt, so die Argumentation, spricht dies für eine Wirklichkeit, in der Moral nicht allein auf menschlichen Einstellungen oder natürlichen Entwicklungsprozessen beruht.
	
	Das Argument verbindet diese Überlegung mit der Annahme, dass ohne Gott der evolutionäre Naturalismus gilt. Daraus ergibt sich die folgende vereinfachte Argumentstruktur:
	
	\begin{tcolorbox}[title=Vereinfachte Struktur des moralischen Arguments]
		\begin{enumerate}[leftmargin=*]
			\item Wenn Gott nicht existiert, gilt der evolutionäre Naturalismus.
			\item Wenn der evolutionäre Naturalismus gilt, existieren keine objektiven moralischen Werte.
			\item Objektive moralische Werte existieren.
			\item Daher existiert Gott.
		\end{enumerate}
	\end{tcolorbox}
	
	Formal ist dieser Schluss gültig, sofern die Aussagen wie angegeben miteinander verknüpft werden. Entscheidend ist jedoch, ob die einzelnen Voraussetzungen wahr sind.
	
	Insbesondere ist die zweite Prämisse umstritten. Sie setzt voraus, dass objektive Moral mit dem evolutionären Naturalismus unvereinbar ist. Eine naturalistische Position könnte dagegen versuchen, objektive moralische Tatsachen unabhängig von Gott zu begründen. Die Frage, ob ein solcher Ansatz erfolgreich ist, muss gesondert untersucht werden.
	
	\section{Die probabilistische Perspektive}
	
	Der ursprüngliche Argumentationsansatz lässt sich zusätzlich mithilfe der Wahrscheinlichkeitstheorie betrachten. Dabei geht es nicht nur darum, ob ein Argument logisch gültig ist, sondern auch darum, wie stark bestimmte Beobachtungen für oder gegen eine Hypothese sprechen.
	
	Man kann beispielsweise fragen, ob die Existenz objektiver Moral unter der Annahme eines Gottes wahrscheinlicher wäre als unter der Annahme, dass kein Gott existiert. Die Antwort hängt von den zugrunde gelegten Wahrscheinlichkeiten ab.
	
	\begin{tcolorbox}[title=Mathematische Erläuterung: Bayes'sche Wahrscheinlichkeiten]
		Der Satz von Bayes beschreibt, wie neue Evidenz eine ursprüngliche Einschätzung verändern kann. In Form von Chancenverhältnissen lautet die Beziehung:
		
		\[
		\frac{P(H\mid E)}{P(\neg H\mid E)}
		=
		\frac{P(H)}{P(\neg H)}
		\cdot
		\frac{P(E\mid H)}{P(E\mid\neg H)}.
		\]
		
		Dabei bezeichnet \(H\) die Hypothese, dass Gott existiert, und \(E\) die betrachtete Evidenz. Der erste Faktor beschreibt die ursprüngliche Einschätzung der Hypothese. Der zweite Faktor vergleicht, wie wahrscheinlich die Evidenz unter der Annahme von Gottes Existenz und unter deren Verneinung ist.
		
		Ist die Evidenz unter einer Hypothese wahrscheinlicher, kann sie diese Hypothese im Vergleich zur Alternative stützen. Die Formel allein entscheidet jedoch nicht, welche Wahrscheinlichkeiten angemessen sind.
	\end{tcolorbox}
	
	Für das moralische Argument ist daher entscheidend, wie die Existenz objektiver Moral bewertet wird und welche Wahrscheinlichkeit man ihr unter verschiedenen Weltanschauungen zuschreibt. Solche Einschätzungen sind nicht allein mathematisch bestimmbar. Sie beruhen auf philosophischen Annahmen, die begründet und kritisch geprüft werden müssen.
	
	Die probabilistische Darstellung kann das Argument strukturieren, liefert aber keinen unabhängigen Beweis für die Existenz Gottes. Insbesondere werden umstrittene Voraussetzungen nicht dadurch gesichert, dass sie in eine mathematische Formel eingesetzt werden.
	
	\section{Zentrale Einwände}
	
	Ein wichtiger Einwand betrifft die Annahme, dass objektive Moral eine göttliche Grundlage benötigt. Selbst wenn moralische Tatsachen objektiv existieren, folgt daraus nicht unmittelbar, dass Gott existieren muss. Es müsste zusätzlich gezeigt werden, dass alternative Erklärungen unzureichend sind.
	
	Ein zweiter Einwand betrifft die evolutionäre Erklärung moralischer Überzeugungen. Dass sich moralische Intuitionen im Laufe der Evolution entwickelt haben, widerlegt ihre objektive Gültigkeit nicht automatisch. Evolutionäre Prozesse könnten beispielsweise Fähigkeiten hervorgebracht haben, die Menschen dabei helfen, moralisch relevante Aspekte ihrer Umwelt zu erkennen. Ob diese Erklärung ausreicht, bleibt offen.
	
	Ein dritter Einwand ist das klassische Euthyphron-Dilemma. Es fragt, ob etwas gut ist, weil Gott es gebietet, oder ob Gott es gebietet, weil es gut ist.
	
	Im ersten Fall scheint Moral von Gottes Willen abzuhängen und könnte dadurch willkürlich wirken. Im zweiten Fall scheint das Gute unabhängig von Gott zu bestehen. Das Dilemma stellt deshalb infrage, ob Gott tatsächlich die notwendige Grundlage objektiver Moral ist.
	
	Theistische Ansätze versuchen, dieses Problem beispielsweise dadurch zu beantworten, dass sie das Gute nicht als willkürlichen Befehl, sondern als Ausdruck von Gottes Wesen verstehen. Ob diese Antwort das Problem löst, ist jedoch ebenfalls umstritten.
	
	Schliesslich ist auch die Ablehnung des moralischen Realismus eine mögliche Reaktion. Wer bestreitet, dass objektive moralische Tatsachen existieren, muss die zentrale Voraussetzung des Arguments nicht akzeptieren. Damit ist das Argument für diese Person nicht überzeugend, ohne dass dadurch sämtliche Fragen der Moral geklärt wären.
	
	\section{Fazit}
	
	Das moralische Argument untersucht, ob die Existenz objektiver moralischer Werte mit dem evolutionären Naturalismus vereinbar ist und ob der Theismus eine überzeugendere Erklärung bietet. Sein Ausgangspunkt ist die Unterscheidung zwischen der evolutionären Entstehung moralischer Überzeugungen und der Frage, ob diese Überzeugungen objektiv wahr sind.
	
	Das Argument kann auf eine philosophische Herausforderung für den Naturalismus hinweisen. Es beweist jedoch nicht, dass objektive Moral unmöglich naturalistisch erklärt werden kann. Ebenso folgt aus der Existenz objektiver Moral nicht ohne weitere Begründung, dass Gott existiert.
	
	Die Stärke des Arguments hängt daher vor allem von zwei Fragen ab: Gibt es objektive moralische Tatsachen, und benötigen diese eine göttliche Grundlage? Solange beide Fragen umstritten bleiben, ist das moralische Argument als philosophische Verteidigung des Theismus zu verstehen, nicht als allgemein anerkannter Beweis für die Existenz Gottes.
	
	\begin{thebibliography}{9}
		
		\bibitem{craig}
		Craig, W. L. (2009).
		\textit{Reasonable Faith: Christian Truth and Apologetics}.
		3. Auflage. Crossway.
		
		\bibitem{joyce}
		Joyce, R. (2006).
		\textit{The Evolution of Morality}.
		MIT Press.
		
		\bibitem{linville}
		Linville, M. D. (2009).
		The Moral Argument.
		In: W. L. Craig und J. P. Moreland (Hrsg.),
		\textit{The Blackwell Companion to Natural Theology},
		S. 391--448.
		
		\bibitem{mawson}
		Mawson, T. J. (2009).
		Morality and Religion.
		\textit{Philosophy Compass}, 4(6), 1033--1043.
		
		\bibitem{sober}
		Sober, E. (1994).
		\textit{From a Biological Point of View}.
		Cambridge University Press.
		
	\end{thebibliography}
	
\end{document}